% !TeX root = ../proposal.tex
% Rules 6.2, 7.6, and 21: List I/O and request dedicated signal pads within the three-pad limit.
% Rules 7.5/22: Up to four shared analog buses; request one for the buffered IF output.
% RF uses two dedicated pads; single-ended LO uses one pad and an on-chip balun.
% Asymmetric external interfaces are compatible with a double-balanced internal core.
% Direct differential RF avoids an on-chip conversion stage ahead of the RF transconductor.
% A single-ended bench source still needs an external RF balun or differential fixture.
% Rule 7.4 provides digital control resources; ENABLE does not consume an analog signal pad.
% Rules 7.1/10 specify 1.8 V digital and 3.3 V analog harness domains and I/O protection.
% A 1.8 V analog core is our design choice; its actual supply connection must be confirmed.
% Internal bias references use Rules 7.2/7.3, subject to final allocation.

\section{Pinout and I/O Requirements}

We request three dedicated analog signal pads: two for differential RF and one for single-ended LO, plus one shared analog signal connection for the buffered IF output.

\begin{table}[htbp]
\centering
\caption{Proposed I/O signals and interfaces}
\label{tab:io}
\begin{tabularx}{\textwidth}{@{}llY@{}}
\toprule
\textbf{Signal} &
\textbf{Direction} &
\textbf{Description} \\
\midrule

\texttt{RF\_P / RF\_N} &
Input &
Differential RF input, nominally \(900\,\mathrm{MHz}\); two dedicated pads, no on-chip RF balun \\

\texttt{LO\_IN} &
Input &
Single-ended external LO input with on-chip balun and buffer \\

\texttt{IF\_OUT} &
Output &
Buffered, single-ended IF output through one shared analog bus; nominally \(10\,\mathrm{MHz}\), evaluated over \(5\)--\(20\,\mathrm{MHz}\) \\

\texttt{VDD / GND} &
Power &
Supplied through the Chipalooza harness \\

\texttt{ENABLE} &
Digital input &
Mixer enable/power-down control, if available \\

\texttt{VREF / IREF} &
Analog reference &
Harness bias references; final allocation to be confirmed \\

\bottomrule
\end{tabularx}
\end{table}

RF is supplied differentially, while the on-chip LO balun generates differential LO drive.
Both internal signal paths remain differential, as required by the double-balanced core; their external interfaces need not be identical.
Using the extra pad for RF removes the on-chip RF conversion stage and its noise, distortion, and power contributions, while retaining single-ended LO drive.

\textbf{Pad requirements:}
\texttt{RF\_P}, \texttt{RF\_N}, and \texttt{LO\_IN} require dedicated connections suitable for \(900\,\mathrm{MHz}\) signaling.
\texttt{IF\_OUT} is proposed to use one shared analog bus, with a buffer designed for the complete output path and specified external load.
The shared path will be verified for gain error, phase response, noise, and distortion over \(5\)--\(20\,\mathrm{MHz}\), including switch resistance, interconnect and pad capacitance, package parasitics, and output-buffer impedance.
All three dedicated pads are used, so a dedicated IF fallback would require revising the RF interface to single-ended with an on-chip RF balun, or obtaining an approved resource exception.
The shared-path feasibility must therefore be resolved before finalizing this allocation.

The baseline circuit will target a \(1.8\,\mathrm{V}\) core supply, with input and output protection compatible with the harness pad interface.
External RF and LO signal generators will be used for characterization.
An on-chip PLL or LO frequency synthesizer is outside the scope of this proposal.
